\documentclass[UTF8]{ctexart}

\usepackage{amsmath,amssymb}
\usepackage{geometry}
\usepackage{booktabs}
\usepackage{array}
\usepackage{enumitem}
\usepackage[most]{tcolorbox}

\geometry{a4paper, margin=2.2cm}
\setlist[itemize]{itemsep=2pt, topsep=2pt}
\setlist[enumerate]{itemsep=2pt, topsep=2pt}

\title{信号与系统知识点整理}
\author{}
\date{}

\begin{document}

\maketitle
\tableofcontents

\newpage

%========================
\section{常用信号与基本概念}

\subsection{常见连续时间信号}

信号与系统中常见的基本信号有：

\begin{tcolorbox}[colback=blue!5!white,colframe=blue!60!black,title=常见信号]
\begin{itemize}
    \item 正弦信号：$x(t)=A\sin(\omega t+\varphi)$ 或 $x(t)=A\cos(\omega t+\varphi)$
    \item 指数信号：$x(t)=e^{-at}u(t)$
    \item 单位阶跃信号：$u(t)$ 或 $\varepsilon(t)$
    \item 单位冲激信号：$\delta(t)$
    \item 抽样函数：
    \[
    Sa(t)=\frac{\sin t}{t}
    \]
\end{itemize}
\end{tcolorbox}

其中，正弦信号的周期为：

\[
T=\frac{2\pi}{\omega}
\]

若两个周期信号相加，且两个周期之比为有理数，则其和仍然是周期信号，周期为两个周期的最小公倍周期。

例如：

\[
x_1(t)=\sin \omega_1 t,\qquad x_2(t)=\cos \omega_2 t
\]

若：

\[
\frac{T_1}{T_2} \in \mathbb{Q}
\]

则：

\[
x(t)=x_1(t)+x_2(t)
\]

为周期信号。

\subsection{单位阶跃信号}

单位阶跃信号定义为：

\[
u(t)=
\begin{cases}
0, & t<0\\
1, & t>0
\end{cases}
\]

常用形式：

\[
u(t-a)=
\begin{cases}
0, & t<a\\
1, & t>a
\end{cases}
\]

表示信号在 $t=a$ 时刻开始。

\[
u(a-t)=
\begin{cases}
1, & t<a\\
0, & t>a
\end{cases}
\]

表示信号在 $t=a$ 时刻结束。

\subsection{矩形窗表示法}

利用阶跃函数可以表示有限时间区间内的信号。

若信号只在区间 $a<t<b$ 内存在，则可以写成：

\[
x(t)=f(t)\left[u(t-a)-u(t-b)\right]
\]

例如，信号 $f(t)$ 只在 $0<t<3$ 内存在：

\[
x(t)=f(t)\left[u(t)-u(t-3)\right]
\]

如果信号在 $-5<t<0$ 内存在，则：

\[
x(t)=f(t)\left[u(t+5)-u(t)\right]
\]

\subsection{单位冲激信号}

单位冲激信号 $\delta(t)$ 具有如下性质：

\begin{tcolorbox}[colback=yellow!8!white,colframe=orange!80!black,title=冲激信号的重要性质]
\[
\int_{-\infty}^{+\infty}\delta(t)\,dt=1
\]

\[
\int_{-\infty}^{+\infty}f(t)\delta(t)\,dt=f(0)
\]

\[
\int_{-\infty}^{+\infty}f(t)\delta(t-t_0)\,dt=f(t_0)
\]
\end{tcolorbox}

这叫做冲激函数的筛选性质，也叫抽样性质。

\subsection{冲激函数与阶跃函数的关系}

单位冲激函数是单位阶跃函数的导数：

\[
\delta(t)=\frac{d}{dt}u(t)
\]

单位阶跃函数是单位冲激函数的积分：

\[
u(t)=\int_{-\infty}^{t}\delta(\tau)\,d\tau
\]

对于移位形式：

\[
\frac{d}{dt}u(t-a)=\delta(t-a)
\]

%========================
\section{信号的基本运算与波形变换}

\subsection{信号的平移}

若原信号为 $f(t)$，则：

\[
f(t-t_0)
\]

表示信号向右平移 $t_0$。

\[
f(t+t_0)
\]

表示信号向左平移 $t_0$。

\begin{tcolorbox}[colback=green!5!white,colframe=green!50!black,title=平移判断]
\begin{itemize}
    \item $f(t-a)$：向右移 $a$
    \item $f(t+a)$：向左移 $a$
\end{itemize}
\end{tcolorbox}

\subsection{信号的反转}

\[
f(-t)
\]

表示信号关于纵轴反转。

也就是说，原来在 $t=a$ 处的特征，会变到 $t=-a$ 处。

\subsection{尺度变换}

\[
f(at)
\]

表示时间尺度变化。

\begin{itemize}
    \item 当 $|a|>1$ 时，波形在时间轴上压缩。
    \item 当 $0<|a|<1$ 时，波形在时间轴上展宽。
    \item 当 $a<0$ 时，还包含关于纵轴反转。
\end{itemize}

例如：

\[
f(2t)
\]

表示信号压缩为原来的一半。

\[
f\left(\frac{t}{2}\right)
\]

表示信号展宽为原来的两倍。

\subsection{波形变换的一般步骤}

处理复杂形式：

\[
f(at+b)
\]

时，可以按如下步骤：

\begin{enumerate}
    \item 先将括号内整理成：
    \[
    f\left[a\left(t+\frac{b}{a}\right)\right]
    \]
    \item 判断是否反转：看 $a$ 的正负。
    \item 判断压缩或展宽：看 $|a|$。
    \item 判断平移：看 $t+\frac{b}{a}$。
\end{enumerate}

\subsection{波形微分}

若信号存在突变，则微分后会出现冲激信号。

例如：

\[
\frac{d}{dt}u(t-a)=\delta(t-a)
\]

若矩形信号为：

\[
x(t)=u(t)-u(t-T)
\]

则其导数为：

\[
x'(t)=\delta(t)-\delta(t-T)
\]

也就是说，矩形波上升沿对应正冲激，下降沿对应负冲激。

%========================
\section{系统响应}

\subsection{系统响应的分类}

线性系统的完全响应通常可分为：

\[
y(t)=y_{zi}(t)+y_{zs}(t)
\]

其中：

\begin{itemize}
    \item $y_{zi}(t)$：零输入响应，由系统初始状态引起。
    \item $y_{zs}(t)$：零状态响应，由外加激励引起。
\end{itemize}

也可以从响应形态分为：

\[
y(t)=y_n(t)+y_f(t)
\]

其中：

\begin{itemize}
    \item $y_n(t)$：自然响应，与系统固有特性有关。
    \item $y_f(t)$：强迫响应，与输入信号有关。
\end{itemize}

还可以从时间变化趋势分为：

\begin{itemize}
    \item 瞬态响应：随时间逐渐衰减的部分。
    \item 稳态响应：系统稳定后保留下来的部分。
\end{itemize}

\subsection{线性时不变系统}

线性时不变系统简称 LTI 系统。

其性质包括：

\begin{itemize}
    \item 线性：满足叠加性和齐次性。
    \item 时不变：输入延时，输出也相同延时。
\end{itemize}

若系统输入为 $f(t)$，单位冲激响应为 $h(t)$，则零状态响应为：

\[
y_{zs}(t)=f(t)*h(t)
\]

其中 $*$ 表示卷积。

\subsection{卷积积分}

连续时间卷积定义为：

\[
y(t)=f(t)*h(t)=\int_{-\infty}^{+\infty}f(\tau)h(t-\tau)\,d\tau
\]

卷积的基本性质：

\begin{tcolorbox}[colback=blue!5!white,colframe=blue!60!black,title=卷积性质]
\[
f(t)*h(t)=h(t)*f(t)
\]

\[
f(t)*\delta(t)=f(t)
\]

\[
f(t)*\delta(t-t_0)=f(t-t_0)
\]

\[
f(t-t_1)*\delta(t-t_2)=f(t-t_1-t_2)
\]
\end{tcolorbox}

\subsection{系统连接}

若两个系统的冲激响应分别为 $h_1(t)$ 和 $h_2(t)$：

\begin{itemize}
    \item 串联系统：
    \[
    h(t)=h_1(t)*h_2(t)
    \]
    \item 并联系统：
    \[
    h(t)=h_1(t)+h_2(t)
    \]
\end{itemize}

%========================
\section{傅里叶级数与傅里叶变换}

\subsection{周期信号的判断}

若信号 $x(t)$ 满足：

\[
x(t+T)=x(t)
\]

则称 $x(t)$ 为周期信号，其中 $T$ 为周期。

正弦信号：

\[
x(t)=\sin \omega_0 t
\]

的周期为：

\[
T=\frac{2\pi}{\omega_0}
\]

对于多个周期信号相加，若各周期的比值为有理数，则总信号为周期信号。

\subsection{奇偶性}

信号的奇偶性：

\[
x(-t)=x(t)
\]

则 $x(t)$ 为偶函数。

\[
x(-t)=-x(t)
\]

则 $x(t)$ 为奇函数。

常见结论：

\begin{itemize}
    \item $\cos \omega t$ 是偶函数。
    \item $\sin \omega t$ 是奇函数。
    \item 偶函数的傅里叶级数只含余弦项。
    \item 奇函数的傅里叶级数只含正弦项。
\end{itemize}

\subsection{三角形式傅里叶级数}

周期信号 $x(t)$ 可以展开为：

\[
x(t)=a_0+\sum_{n=1}^{\infty}\left[a_n\cos(n\omega_0 t)+b_n\sin(n\omega_0 t)\right]
\]

其中：

\[
\omega_0=\frac{2\pi}{T}
\]

系数为：

\[
a_0=\frac{1}{T}\int_{T}x(t)\,dt
\]

\[
a_n=\frac{2}{T}\int_{T}x(t)\cos(n\omega_0 t)\,dt
\]

\[
b_n=\frac{2}{T}\int_{T}x(t)\sin(n\omega_0 t)\,dt
\]

若 $x(t)$ 为偶函数，则：

\[
b_n=0
\]

若 $x(t)$ 为奇函数，则：

\[
a_0=0,\qquad a_n=0
\]

\subsection{傅里叶变换定义}

连续时间傅里叶变换定义为：

\[
F(\omega)=\int_{-\infty}^{+\infty}f(t)e^{-j\omega t}\,dt
\]

反变换为：

\[
f(t)=\frac{1}{2\pi}\int_{-\infty}^{+\infty}F(\omega)e^{j\omega t}\,d\omega
\]

记作：

\[
f(t)\leftrightarrow F(\omega)
\]

\subsection{常用傅里叶变换对}

\begin{table}[h]
\centering
\begin{tabular}{c c}
\toprule
时域信号 & 频域信号 \\
\midrule
$\delta(t)$ & $1$ \\
$1$ & $2\pi\delta(\omega)$ \\
$e^{-at}u(t),\ a>0$ & $\dfrac{1}{a+j\omega}$ \\
$u(t)$ & $\pi\delta(\omega)+\dfrac{1}{j\omega}$ \\
$\dfrac{\sin Wt}{\pi t}$ & 矩形频谱 \\
$Sa(Wt)=\dfrac{\sin Wt}{Wt}$ & 矩形频谱形式 \\
\bottomrule
\end{tabular}
\end{table}

\subsection{傅里叶变换性质}

\subsubsection{线性性质}

若：

\[
f_1(t)\leftrightarrow F_1(\omega)
\]

\[
f_2(t)\leftrightarrow F_2(\omega)
\]

则：

\[
af_1(t)+bf_2(t)\leftrightarrow aF_1(\omega)+bF_2(\omega)
\]

\subsubsection{时移性质}

若：

\[
f(t)\leftrightarrow F(\omega)
\]

则：

\[
f(t-t_0)\leftrightarrow F(\omega)e^{-j\omega t_0}
\]

时域平移对应频域乘以相位因子。

\subsubsection{频移性质}

若：

\[
f(t)\leftrightarrow F(\omega)
\]

则：

\[
f(t)e^{j\omega_0 t}\leftrightarrow F(\omega-\omega_0)
\]

\[
f(t)e^{-j\omega_0 t}\leftrightarrow F(\omega+\omega_0)
\]

\subsubsection{尺度变换}

若：

\[
f(t)\leftrightarrow F(\omega)
\]

则：

\[
f(at)\leftrightarrow \frac{1}{|a|}F\left(\frac{\omega}{a}\right)
\]

时域压缩，频域展宽；时域展宽，频域压缩。

\subsubsection{微分性质}

若：

\[
f(t)\leftrightarrow F(\omega)
\]

则：

\[
\frac{d f(t)}{dt}\leftrightarrow j\omega F(\omega)
\]

\subsubsection{积分性质}

若：

\[
f(t)\leftrightarrow F(\omega)
\]

则：

\[
\int_{-\infty}^{t}f(\tau)\,d\tau
\leftrightarrow
\frac{F(\omega)}{j\omega}+\pi F(0)\delta(\omega)
\]

\subsection{调制性质}

由于：

\[
\cos \omega_0 t=\frac{1}{2}\left(e^{j\omega_0 t}+e^{-j\omega_0 t}\right)
\]

所以：

\[
f(t)\cos \omega_0 t
\leftrightarrow
\frac{1}{2}\left[F(\omega-\omega_0)+F(\omega+\omega_0)\right]
\]

由于：

\[
\sin \omega_0 t=\frac{1}{2j}\left(e^{j\omega_0 t}-e^{-j\omega_0 t}\right)
\]

所以：

\[
f(t)\sin \omega_0 t
\leftrightarrow
\frac{1}{2j}\left[F(\omega-\omega_0)-F(\omega+\omega_0)\right]
\]

\subsection{采样定理}

若连续信号 $f(t)$ 的最高频率为 $f_{\max}$，为了能够由采样信号无失真恢复原信号，采样频率必须满足：

\[
f_s\geq 2f_{\max}
\]

其中：

\[
f_s=\frac{1}{T_s}
\]

称 $2f_{\max}$ 为奈奎斯特采样频率。

%========================
\section{拉普拉斯变换}

\subsection{拉普拉斯变换定义}

连续信号 $f(t)$ 的单边拉普拉斯变换为：

\[
F(s)=\int_{0^-}^{+\infty}f(t)e^{-st}\,dt
\]

记作：

\[
f(t)\leftrightarrow F(s)
\]

\subsection{常用拉普拉斯变换对}

\begin{table}[h]
\centering
\begin{tabular}{c c}
\toprule
时域信号 & 拉普拉斯变换 \\
\midrule
$\delta(t)$ & $1$ \\
$u(t)$ & $\dfrac{1}{s}$ \\
$e^{-at}u(t)$ & $\dfrac{1}{s+a}$ \\
$t u(t)$ & $\dfrac{1}{s^2}$ \\
$\dfrac{t^{n-1}}{(n-1)!}u(t)$ & $\dfrac{1}{s^n}$ \\
$\sin \omega_0 t\,u(t)$ & $\dfrac{\omega_0}{s^2+\omega_0^2}$ \\
$\cos \omega_0 t\,u(t)$ & $\dfrac{s}{s^2+\omega_0^2}$ \\
\bottomrule
\end{tabular}
\end{table}

\subsection{拉普拉斯变换性质}

\subsubsection{线性性质}

\[
af_1(t)+bf_2(t)\leftrightarrow aF_1(s)+bF_2(s)
\]

\subsubsection{时移性质}

\[
f(t-t_0)u(t-t_0)\leftrightarrow e^{-st_0}F(s)
\]

\subsubsection{复频移性质}

若：

\[
f(t)\leftrightarrow F(s)
\]

则：

\[
e^{-at}f(t)\leftrightarrow F(s+a)
\]

\subsubsection{尺度性质}

\[
f(at)\leftrightarrow \frac{1}{a}F\left(\frac{s}{a}\right),\qquad a>0
\]

\subsubsection{微分性质}

单边拉普拉斯变换中：

\[
f'(t)\leftrightarrow sF(s)-f(0^-)
\]

\[
f''(t)\leftrightarrow s^2F(s)-sf(0^-)-f'(0^-)
\]

一般地：

\[
f^{(n)}(t)\leftrightarrow s^nF(s)-s^{n-1}f(0^-)-s^{n-2}f'(0^-)-\cdots-f^{(n-1)}(0^-)
\]

\subsection{由系统函数求微分方程}

若系统函数为：

\[
H(s)=\frac{Y(s)}{F(s)}
=\frac{b_ms^m+b_{m-1}s^{m-1}+\cdots+b_0}
{a_ns^n+a_{n-1}s^{n-1}+\cdots+a_0}
\]

则有：

\[
\left(a_ns^n+a_{n-1}s^{n-1}+\cdots+a_0\right)Y(s)
=
\left(b_ms^m+b_{m-1}s^{m-1}+\cdots+b_0\right)F(s)
\]

将 $s$ 替换成微分算子，即可得到微分方程：

\[
a_ny^{(n)}(t)+a_{n-1}y^{(n-1)}(t)+\cdots+a_0y(t)
=
b_mf^{(m)}(t)+b_{m-1}f^{(m-1)}(t)+\cdots+b_0f(t)
\]

\subsection{$s$ 域电路模型}

在拉普拉斯变换中，电路元件可以等效为 $s$ 域阻抗：

\begin{table}[h]
\centering
\begin{tabular}{c c}
\toprule
元件 & $s$ 域阻抗 \\
\midrule
电阻 $R$ & $R$ \\
电容 $C$ & $\dfrac{1}{sC}$ \\
电感 $L$ & $sL$ \\
\bottomrule
\end{tabular}
\end{table}

因此，求电路响应时可以按如下步骤：

\begin{enumerate}
    \item 将电路变换到 $s$ 域。
    \item 电阻仍为 $R$。
    \item 电容用 $\dfrac{1}{sC}$ 表示。
    \item 电感用 $sL$ 表示。
    \item 根据串并联关系列写电压或电流方程。
    \item 求出 $Y(s)$ 后，再进行拉普拉斯反变换。
\end{enumerate}

\subsection{系统稳定性}

连续时间 LTI 系统稳定的充要条件是：

\[
H(s)
\]

的所有极点均位于 $s$ 平面的左半平面，即：

\[
\operatorname{Re}(s_i)<0
\]

如果存在极点位于右半平面，则系统不稳定。

如果极点位于虚轴上，需要具体分析，通常不能简单判为稳定。

%========================
\section{离散时间信号与系统}

\subsection{常见离散信号}

离散时间信号通常写作：

\[
f(k)
\]

其中 $k$ 为整数。

常见离散信号有：

\begin{itemize}
    \item 单位序列：
    \[
    \delta(k)=
    \begin{cases}
    1, & k=0\\
    0, & k\neq 0
    \end{cases}
    \]
    \item 单位阶跃序列：
    \[
    u(k)=
    \begin{cases}
    1, & k\geq 0\\
    0, & k<0
    \end{cases}
    \]
    \item 指数序列：
    \[
    a^ku(k)
    \]
\end{itemize}

\subsection{单位序列的筛选性质}

离散单位序列具有筛选性质：

\[
f(k)\delta(k-k_0)=f(k_0)\delta(k-k_0)
\]

并且：

\[
\sum_{k=-\infty}^{+\infty}f(k)\delta(k-k_0)=f(k_0)
\]

\subsection{离散卷积和}

离散时间 LTI 系统的零状态响应为：

\[
y_{zs}(k)=f(k)*h(k)
\]

其中卷积和定义为：

\[
y(k)=\sum_{m=-\infty}^{+\infty}f(m)h(k-m)
\]

常用性质：

\[
f(k)*\delta(k)=f(k)
\]

\[
f(k)*\delta(k-k_0)=f(k-k_0)
\]

\[
f(k-k_1)*\delta(k-k_2)=f(k-k_1-k_2)
\]

\subsection{离散系统差分方程}

离散系统常用差分方程描述：

\[
y(k)+a_1y(k-1)+a_2y(k-2)+\cdots+a_Ny(k-N)
=
b_0f(k)+b_1f(k-1)+\cdots+b_Mf(k-M)
\]

其中：

\begin{itemize}
    \item $y(k-i)$ 表示输出延时 $i$ 个单位。
    \item $f(k-i)$ 表示输入延时 $i$ 个单位。
    \item 延时单元在 $z$ 域中对应 $z^{-1}$。
\end{itemize}

%========================
\section{$Z$ 变换}

\subsection{$Z$ 变换定义}

离散序列 $f(k)$ 的双边 $Z$ 变换定义为：

\[
F(z)=\sum_{k=-\infty}^{+\infty}f(k)z^{-k}
\]

记作：

\[
f(k)\leftrightarrow F(z)
\]

单边 $Z$ 变换为：

\[
F(z)=\sum_{k=0}^{+\infty}f(k)z^{-k}
\]

\subsection{常用 $Z$ 变换对}

\begin{table}[h]
\centering
\begin{tabular}{c c}
\toprule
时域序列 & $Z$ 变换 \\
\midrule
$\delta(k)$ & $1$ \\
$u(k)$ & $\dfrac{1}{1-z^{-1}}$ \\
$a^ku(k)$ & $\dfrac{1}{1-az^{-1}}$ \\
$-a^ku(-k-1)$ & $\dfrac{1}{1-az^{-1}}$ \\
\bottomrule
\end{tabular}
\end{table}

需要注意：同一个代数表达式可能对应不同序列，关键取决于收敛域。

\subsection{$Z$ 变换时移性质}

若：

\[
f(k)\leftrightarrow F(z)
\]

则双边 $Z$ 变换中有：

\[
f(k-k_0)\leftrightarrow z^{-k_0}F(z)
\]

特别地：

\[
f(k-1)\leftrightarrow z^{-1}F(z)
\]

\[
f(k+1)\leftrightarrow zF(z)
\]

\subsection{收敛域}

$Z$ 变换必须结合收敛域 ROC 判断原序列。

例如：

\[
F(z)=\frac{1}{1-az^{-1}}
\]

若收敛域为：

\[
|z|>|a|
\]

则对应右边序列：

\[
f(k)=a^ku(k)
\]

若收敛域为：

\[
|z|<|a|
\]

则对应左边序列：

\[
f(k)=-a^ku(-k-1)
\]

\subsection{利用 $Z$ 变换求差分方程响应}

对于差分方程：

\[
y(k)+a_1y(k-1)+\cdots+a_Ny(k-N)
=
b_0f(k)+b_1f(k-1)+\cdots+b_Mf(k-M)
\]

在零初始条件下进行 $Z$ 变换：

\[
Y(z)\left(1+a_1z^{-1}+\cdots+a_Nz^{-N}\right)
=
F(z)\left(b_0+b_1z^{-1}+\cdots+b_Mz^{-M}\right)
\]

所以系统函数为：

\[
H(z)=\frac{Y(z)}{F(z)}
=
\frac{b_0+b_1z^{-1}+\cdots+b_Mz^{-M}}
{1+a_1z^{-1}+\cdots+a_Nz^{-N}}
\]

%========================
\section{系统函数与稳定性}

\subsection{连续系统函数}

连续系统中：

\[
H(s)=\frac{Y(s)}{F(s)}
\]

称为系统函数。

若输入为 $f(t)$，输出为 $y(t)$，则：

\[
Y(s)=H(s)F(s)
\]

\subsection{离散系统函数}

离散系统中：

\[
H(z)=\frac{Y(z)}{F(z)}
\]

若输入为 $f(k)$，输出为 $y(k)$，则：

\[
Y(z)=H(z)F(z)
\]

\subsection{连续系统稳定性}

连续系统稳定条件：

\[
H(s)\text{ 的所有极点均在左半平面}
\]

即：

\[
\operatorname{Re}(p_i)<0
\]

\subsection{离散系统稳定性}

离散系统稳定条件：

\[
H(z)\text{ 的所有极点均在单位圆内}
\]

即：

\[
|p_i|<1
\]

%========================
\section{复习重点总结}

\begin{tcolorbox}[colback=red!5!white,colframe=red!60!black,title=重点一：波形变换]
重点掌握：
\[
f(t-a),\quad f(t+a),\quad f(-t),\quad f(at+b)
\]
的平移、反转、尺度变化。
\end{tcolorbox}

\begin{tcolorbox}[colback=red!5!white,colframe=red!60!black,title=重点二：冲激函数]
重点掌握：
\[
\int_{-\infty}^{+\infty}f(t)\delta(t-t_0)\,dt=f(t_0)
\]
以及：
\[
f(t)*\delta(t-t_0)=f(t-t_0)
\]
\end{tcolorbox}

\begin{tcolorbox}[colback=red!5!white,colframe=red!60!black,title=重点三：卷积]
连续卷积：
\[
y(t)=\int_{-\infty}^{+\infty}f(\tau)h(t-\tau)\,d\tau
\]

离散卷积：
\[
y(k)=\sum_{m=-\infty}^{+\infty}f(m)h(k-m)
\]
\end{tcolorbox}

\begin{tcolorbox}[colback=red!5!white,colframe=red!60!black,title=重点四：傅里叶变换]
重点掌握：
\[
f(t-t_0)\leftrightarrow F(\omega)e^{-j\omega t_0}
\]

\[
f(at)\leftrightarrow \frac{1}{|a|}F\left(\frac{\omega}{a}\right)
\]

\[
f(t)\cos\omega_0t
\leftrightarrow
\frac{1}{2}\left[F(\omega-\omega_0)+F(\omega+\omega_0)\right]
\]
\end{tcolorbox}

\begin{tcolorbox}[colback=red!5!white,colframe=red!60!black,title=重点五：拉普拉斯变换]
重点掌握：
\[
\delta(t)\leftrightarrow 1
\]

\[
u(t)\leftrightarrow \frac{1}{s}
\]

\[
e^{-at}u(t)\leftrightarrow \frac{1}{s+a}
\]

\[
f'(t)\leftrightarrow sF(s)-f(0^-)
\]
\end{tcolorbox}

\begin{tcolorbox}[colback=red!5!white,colframe=red!60!black,title=重点六：$Z$ 变换]
重点掌握：
\[
f(k-1)\leftrightarrow z^{-1}F(z)
\]

\[
\delta(k)\leftrightarrow 1
\]

\[
a^ku(k)\leftrightarrow \frac{1}{1-az^{-1}}
\]
以及收敛域对反变换结果的影响。
\end{tcolorbox}

\end{document}